\documentclass[12pt]{article}
\usepackage[margin=1in]{geometry}
\setlength{\parskip}{0.5em}

\title{Formatted Text in Context}
\author{legible-pdf synthesized corpus}
\date{}

\begin{document}
\maketitle

\section{Cloud Atlas}

Among the principal genera of clouds, \textit{nimbus} clouds produce
precipitation, while \textbf{stratus} clouds form low layered sheets across
the sky. The Latin token \texttt{cumulus} appears in many compound names
in meteorological taxonomy.

\section{Greek Letter Reference}

The following table summarizes a few Greek letters paired with formatted
representations.

\begin{table}[h]
\centering
\begin{tabular}{|l|l|l|}
\hline
Plain & Formatted & Note \\
\hline
alpha & \textit{beta} & second letter \\
gamma & \textbf{delta} & fourth letter \\
epsilon & \texttt{zeta} & sixth letter \\
\hline
\end{tabular}
\caption{Greek letter reference.}
\end{table}

\section{Cloud Subtypes}

The classification of clouds extends to many compound forms, some of which
are noted in the supporting reference.\footnote{Among the subtypes,
\textit{altocumulus} clouds occupy the middle altitudes, the high
\textbf{cirrus} clouds form delicate filaments, and the boundary form
\texttt{stratocumulus} marks transitions between layers.}

\section{Glossary}

The document mentions nine distinctive tokens: nimbus, stratus, cumulus,
beta, delta, zeta, altocumulus, cirrus, and stratocumulus.

\end{document}
